% BANKS-SOURCE: pdf=71 print=61 kind=solution-section id=chapter05-numbered-solutions
\begin{bankssolutions}{Solutions to Problems for Chapter 5}

% BANKS-SOURCE: pdf=69 print=59 kind=solution id=5.1
\BanksSolution{5.1}
In $D=4$, put $p^0=\omega_{\mathbf p}=(\mathbf p^2+m^2)^{1/2}$, $p_\mu=(-\omega_{\mathbf p},\mathbf p)$,
so $p\cdot x=-\omega_{\mathbf p}t+\mathbf p\cdot\mathbf x$,
and let $s=\pm\frac12$. With the
normalization used in Appendix C, the general operator-valued solution is
\begin{equation*}
  \psi(x)=\int\frac{\dd^{D-1}\mathbf p}{\sqrt{(2\pi)^{D-1}\,2\omega_{\mathbf p}}}
  \sum_s\left[a(\mathbf p,s)u(p,s)\ee^{+\ii p\cdot x}
  +b^\dagger(\mathbf p,s)v(p,s)\ee^{-\ii p\cdot x}\right].
\end{equation*}
Here $(\slashed p-m)u(p,s)=0$, $(\slashed p+m)v(p,s)=0$, and the two values of
$s$ span the respective solution spaces. The adjoint field is obtained by taking
the Hermitian adjoint and multiplying by $\gamma^0$.

For a proper orthochronous Lorentz transformation, write
$R_W(\Lambda,p)=L^{-1}(\Lambda p)\Lambda L(p)$. The spinors are covariantly
normalized, for example
$\bar u(p,r)u(p,s)=-2m\delta_{rs}$ and
$\bar v(p,r)v(p,s)=2m\delta_{rs}$. Their covariance relation therefore has
no energy factor:
\begin{align*}
  S(\Lambda)u(p,s)&=
    \sum_r u(\Lambda p,r)D^{1/2}_{rs}(R_W),\\
  S(\Lambda)v(p,s)&=
    \sum_r v(\Lambda p,r)D^{1/2}_{rs}(R_W),
\end{align*}
where $S(\Lambda)$ is the $[2,1]\oplus[1,2]$ spinor representation. Comparing
the transformed expansion with
$U(\Lambda)^{-1}\psi(x)U(\Lambda)=S(\Lambda)\psi(\Lambda^{-1}x)$ gives
\begin{align*}
  U(\Lambda)a^\dagger(\mathbf p,s)U(\Lambda)^{-1}
  &=\sqrt{\frac{\omega_{\boldsymbol{\Lambda p}}}{\omega_{\mathbf p}}}
   \sum_rD^{1/2}_{rs}(R_W)a^\dagger(\boldsymbol{\Lambda p},r),\\
 U(\Lambda)b^\dagger(\mathbf p,s)U(\Lambda)^{-1}
  &=\sqrt{\frac{\omega_{\boldsymbol{\Lambda p}}}{\omega_{\mathbf p}}}
   \sum_rD^{1/2}_{rs}(R_W)b^\dagger(\boldsymbol{\Lambda p},r).
\end{align*}
The annihilation laws are the adjoints, with the complex-conjugate Wigner
matrix. The square root occurs once, in the law for the delta-normalized
operators. Indeed, if
\begin{equation*}
  d\mu(p)\equiv\frac{\dd^{D-1}\mathbf p}{\sqrt{(2\pi)^{D-1}\,2\omega_{\mathbf p}}},
\end{equation*}
then invariance of $\dd^{D-1}\mathbf p/\omega_{\mathbf p}$ gives
\begin{equation*}
  d\mu(p)\sqrt{\frac{\omega_{\boldsymbol{\Lambda p}}}{\omega_{\mathbf p}}}
  =d\mu(\Lambda p).
\end{equation*}
After the change of variables $q=\Lambda p$, this identity, together with
the spinor intertwiners above, turns the transformed mode expansion into
$S(\Lambda)\psi(\Lambda^{-1}x)$. Thus there is no second energy factor in the
spinor transformation.

Let $\varepsilon=+1$ for a commutator and $\varepsilon=-1$ for an
anticommutator. The canonical algebra forced by the equal-time relation is
\begin{equation*}
 [a(\mathbf p,s),a^\dagger(\mathbf q,r)]_\varepsilon
 =[b(\mathbf p,s),b^\dagger(\mathbf q,r)]_\varepsilon
 =\delta_{sr}\delta^{D-1}(\mathbf p-\mathbf q),
\end{equation*}
with all other brackets zero. The spin sums then give, for $\varepsilon=-1$,
\begin{equation*}
\begin{aligned}
 \{\psi_\alpha(x),\bar\psi_\beta(y)\}
  &=\left(\ii\slashed\partial_x-m\right)_{\alpha\beta}\Delta(x-y),\\
 \Delta(z)&=\int\frac{\dd^{D-1}\mathbf p}{(2\pi)^{D-1}\,2\omega_{\mathbf p}}
  \left(\ee^{+\ii p\cdot z}-\ee^{-\ii p\cdot z}\right).
\end{aligned}
\end{equation*}
The Pauli--Jordan function and all of its derivatives vanish for spacelike $z$.
The analogous anticommutators of two unbarred fields and of two barred fields
vanish as well, so the field is local in the graded sense.

For the commutator choice, the antiparticle term has the opposite sign and gives
the sum of the two frequency distributions rather than their difference. This
Hadamard distribution has support at spacelike separation, so the commutator of
spinor fields does not vanish there. Equivalently, at equal time the proposed
bosonic relation would require a spatial Fourier kernel with both an even mass
term and a nonzero spin term to vanish away from the origin, which it cannot do.
Hence only the fermionic algebra produces a local Dirac field and it gives the
required anticommutation with its conjugate.

% BANKS-SOURCE: pdf=69 print=59 kind=solution id=5.2
\BanksSolution{5.2}
In $D=4$, in the Weyl basis,
\begin{equation*}
 \gamma_W^0=\sigma_1\otimes1,
 \qquad \gamma_W^i=\ii\sigma_2\otimes\sigma_i.
\end{equation*}
The unitary matrix
\begin{equation*}
 S_D=\frac1{\sqrt2}(1-\ii\sigma_2)\otimes1
\end{equation*}
gives the Dirac basis $\gamma_D^\mu=S_D^\dagger\gamma_W^\mu S_D$:
\begin{equation*}
 \gamma_D^0=\sigma_3\otimes1,
 \qquad \gamma_D^i=\ii\sigma_2\otimes\sigma_i.
\end{equation*}
Since $\sigma_1$ and $\sigma_3$ are real while $\sigma_2$ is imaginary, direct
matrix multiplication also gives
\begin{equation*}
 (\gamma_D^\mu)^\dagger=\gamma_D^0\gamma_D^\mu\gamma_D^0,
 \qquad
 (\gamma_{D\mu})^\dagger=\gamma_D^0\gamma_{D\mu}\gamma_D^0.
\end{equation*}
The second relation is preserved by every unitary change of basis related to the
Weyl basis.

A Majorana basis is obtained from the Dirac basis with the unitary block matrix
\begin{equation*}
 S_M=\frac1{\sqrt2}
 \begin{pmatrix}1&\sigma_2\\ \sigma_2&-1\end{pmatrix},
 \qquad
 \gamma_M^\mu=S_M^\dagger\gamma_D^\mu S_M.
\end{equation*}
Writing the matrices in $2\times2$ blocks gives the explicit result
\begin{align*}
 \gamma_M^0&=\begin{pmatrix}0&\sigma_2\\ \sigma_2&0\end{pmatrix},&
 \gamma_M^1&=\begin{pmatrix}\ii\sigma_3&0\\0&\ii\sigma_3\end{pmatrix},\\
 \gamma_M^2&=\begin{pmatrix}0&-\sigma_2\\ \sigma_2&0\end{pmatrix},&
 \gamma_M^3&=\begin{pmatrix}-\ii\sigma_1&0\\0&-\ii\sigma_1\end{pmatrix}.
\end{align*}
Every entry in these four matrices is imaginary, and direct multiplication gives
$(\gamma_M^0)^2=1$, $(\gamma_M^i)^2=-1$, and the vanishing mixed
anticommutators. Therefore $-\ii\gamma_M^\mu\partial_\mu-m$ has real
coefficients. A Majorana solution can consequently be chosen real, and the
Majorana condition is simply $\psi^*=\psi$. Unitarity follows from
$S_M^\dagger S_M=1$, using $\sigma_2^2=1$.

% BANKS-SOURCE: pdf=69 print=59 kind=solution id=5.3
\BanksSolution{5.3}
In $D=4$, it is useful to write the field as
\begin{equation*}
 \psi(x)=\int\frac{\dd^{D-1}\mathbf p}{\sqrt{(2\pi)^{D-1}\,2\omega_{\mathbf p}}}
 \sum_s\left[a(\mathbf p,s)u(p,s)\ee^{+\ii p\cdot x}
 +b^\dagger(\mathbf p,s)v(p,s)\ee^{-\ii p\cdot x}\right].
\end{equation*}
The standard spin-$\frac12$ time-reversal matrix in the Weyl basis is
\begin{equation*}
 \Theta_T=\begin{pmatrix}-\sigma_2&0\\0&-\sigma_2\end{pmatrix},
 \qquad
 \Theta_T^{-1}\gamma^{\mu *}\Theta_T
 =P_T{}^\mu{}_{\nu}\gamma^\nu,
 \qquad P_T=\operatorname{diag}(1,-1,-1,-1).
\end{equation*}
Write $p_T=(p^0,-\mathbf p)$ for the positive-energy momentum with reversed
spatial part. Choose the phases of the positive- and negative-frequency spinors so that,
with $\xi_s=(-1)^{1/2-s}$,
\begin{equation*}
 u(p,s)^*=\xi_s\Theta_Tu(p_T,-s),
 \qquad
 v(p,s)^*=\xi_s\Theta_Tv(p_T,-s).
\end{equation*}
Write the homogeneous field law with its own coefficient as
\begin{equation*}
 T^\dagger\psi(t,\mathbf x)T
 =\zeta_T\Theta_T\psi(-t,\mathbf x).
\end{equation*}
The anti-unitarity of $T$ complex-conjugates numerical coefficients in the mode
expansion. Comparing the two frequency parts gives, in a consistent spin basis,
\begin{align*}
 T^\dagger a^\dagger(\mathbf p,s)T&=\eta_T\xi_sa^\dagger(-\mathbf p,-s),\\
 T^\dagger b^\dagger(\mathbf p,s)T&=\bar\eta_T\xi_sb^\dagger(-\mathbf p,-s),
\end{align*}
with
\begin{equation*}
 \eta_T=\zeta_T^*,
 \qquad
 \bar\eta_T=\zeta_T=\eta_T^*.
\end{equation*}
The complex conjugation in the last equality is forced by anti-unitarity when
$\eta_T$ and $\bar\eta_T$ label creation operators. If a phase is attached to
the field law instead, it is $\zeta_T$; multiplying $\Theta_T$ by a common phase
just relabels $\zeta_T$ and leaves the relation unchanged. Taking the adjoint
of either creation-operator law gives the corresponding annihilation law, with
the phase conjugated.

Parity is represented by
\begin{equation*}
 P^{-1}\psi(t,\mathbf x)P=\eta_P\gamma^0\psi(t,-\mathbf x).
\end{equation*}
The upper and lower components of a rest-frame $v$ spinor have the opposite
$\gamma^0$ eigenvalue from the corresponding $u$ spinor. Consequently
\begin{equation*}
 P^{-1}a^\dagger(\mathbf p,s)P=\eta_Pa^\dagger(-\mathbf p,s),
 \qquad
 P^{-1}b^\dagger(\mathbf p,s)P=\bar\eta_Pb^\dagger(-\mathbf p,s),
 \qquad
 \bar\eta_P=-\eta_P.
\end{equation*}
This is the familiar opposite intrinsic parity of a particle and its
antiparticle.

Let $\mathsf C=\ii\gamma^2\gamma^0$ in the Weyl basis, so that
$\mathsf C^{-1}\gamma^\mu\mathsf C=-(\gamma^\mu)^T$. Charge conjugation has
the homogeneous field law
\begin{equation*}
 C^{-1}\psi(x)C=\eta_C\,\mathsf C\,\bar\psi(x)^T,
\end{equation*}
and therefore
\begin{equation*}
  C^{-1}a^\dagger(\mathbf p,s)C=\eta_C b^\dagger(\mathbf p,s),
 \qquad
  C^{-1}b^\dagger(\mathbf p,s)C=\eta_C^*a^\dagger(\mathbf p,s).
\end{equation*}
The phase can be removed by rephasing $b$, subject to $C^2=1$. In two-component
notation, with $\psi=(\nu_a,\bar\chi^{\dot a})^T$,
\begin{equation*}
 C^{-1}\nu_aC=\eta_C\chi_a,
 \qquad
 C^{-1}\chi_aC=\eta_C^*\nu_a,
\end{equation*}
where the internal representation of $\chi$ is the conjugate of that of $\nu$.
Thus the Dirac-field laws in the Weyl representation are
\begin{align*}
 P^{-1}\psi(t,\mathbf x)P&=\eta_P\gamma^0\psi(t,-\mathbf x),\\
 T^{-1}\psi(t,\mathbf x)T&=\zeta_T\Theta_T\psi(-t,\mathbf x),\\
 C^{-1}\psi(x)C&=\eta_C\mathsf C\bar\psi(x)^T,
\end{align*}
with the anti-particle phase relations above.

% BANKS-SOURCE: pdf=69 print=59 kind=solution id=5.4
\BanksSolution{5.4}
In $D=4$, with $\epsilon^{0123}=+1$, use antisymmetrization with unit weight,
$\gamma^{\mu_1\ldots\mu_k}=\gamma^{[\mu_1}\cdots\gamma^{\mu_k]}$. The Clifford
relations and $\gamma_5=\ii\gamma^0\gamma^1\gamma^2\gamma^3$ give
\begin{equation*}
 \gamma^{\mu\nu\lambda}=-\ii\epsilon^{\mu\nu\lambda\kappa}\gamma_\kappa\gamma_5,
 \qquad
 \gamma^{\mu\nu\lambda\kappa}=-\ii\epsilon^{\mu\nu\lambda\kappa}\gamma_5,
\end{equation*}
with the epsilon convention of Appendix C. Thus the sixteen matrices split into
the scalar, vector, tensor, axial-vector, and pseudoscalar Lorentz types.

For the charge-conjugation matrix $\mathsf C$,
$\mathsf C^{-1}\gamma^\mu\mathsf C=-(\gamma^\mu)^T$. Reversing the order of
$k$ matrices in a bilinear gives
\begin{equation*}
 C^{-1}\!\left(\bar\psi\gamma^{\mu_1\ldots\mu_k}\psi\right)C
 =(-1)^{k(k+1)/2}\bar\psi\gamma^{\mu_1\ldots\mu_k}\psi.
\end{equation*}
The $C$ parities for $k=0,1,2,3,4$ are consequently
$+,-,-,+,+$.

Set $P_\mu{}^\nu=T_\mu{}^\nu=\operatorname{diag}(1,-1,-1,-1)$, with $P$
acting at $(t,\mathbf x)\mapsto(t,-\mathbf x)$ and $T$ at
$(t,\mathbf x)\mapsto(-t,\mathbf x)$. The homogeneous field laws imply
\begin{align*}
 P^{-1}B^{\mu_1\ldots\mu_k}(x)P
 &=P^{\mu_1}{}_{\nu_1}\cdots P^{\mu_k}{}_{\nu_k}
   B^{\nu_1\ldots\nu_k}(t,-\mathbf x),\\
 T^{-1}B^{\mu_1\ldots\mu_k}(x)T
 &=T^{\mu_1}{}_{\nu_1}\cdots T^{\mu_k}{}_{\nu_k}
   B^{\nu_1\ldots\nu_k}(-t,\mathbf x),
\end{align*}
where $B^{\mu_1\ldots\mu_k}=\bar\psi\gamma^{\mu_1\ldots\mu_k}\psi$.
The second formula includes the complex conjugation of coefficients by the
anti-unitary $T$; it is the reason that the raw $\gamma_5$ bilinear is $T$ even,
whereas the Hermitian combination $\ii\bar\psi\gamma_5\psi$ is $T$ odd.
In components this says
\begin{center}
\begin{tabular}{c|c|c|c}
 $B$ & $C$ & $P$ & $T$ \\ \hline
 $\bar\psi\psi$ & $+$ & $+$ & $+$ \\
 $\bar\psi\gamma_5\psi$ & $+$ & $-$ & $+$ \\
 $\bar\psi\gamma^\mu\psi$ & $-$ & $(+,-)$ & $(+,-)$ \\
 $\bar\psi\gamma^\mu\gamma_5\psi$ & $+$ & $(-,+)$ & $(+,-)$ \\
 $\bar\psi\gamma^{\mu\nu}\psi$ & $-$ & $(-,+)$ & $(-,+)$
\end{tabular}
\end{center}
Here each ordered pair lists the time-space components first. For the rank-three
and rank-four products the same component rule applies after using the displayed
epsilon identities; their extra pseudotensor character is carried by $\gamma_5$.

% BANKS-SOURCE: pdf=69 print=59 kind=solution id=5.5
\BanksSolution{5.5}
For a nonzero charge $q$, collect the fields of that charge into a column
$\nu(q)$ and those of charge $-q$ into $\nu(-q)$. The mass term is
\begin{equation*}
 \mathcal L_q=\nu(q)^T\epsilon\,m(q)\nu(-q)+\mathrm{h.c.}
\end{equation*}
Independent unitary changes of basis,
$\nu(q)\mapsto U_q\nu(q)$ and
$\nu(-q)\mapsto U_{-q}\nu(-q)$, send
\begin{equation*}
 m(q)\longmapsto U_q^T m(q)U_{-q}.
\end{equation*}
Take a singular-value decomposition $m(q)=V D W^\dagger$, with $D$ diagonal
and nonnegative. Choosing $U_q=V^*$ and $U_{-q}=W$ gives
\begin{equation*}
 U_q^Tm(q)U_{-q}=D=\operatorname{diag}(m_1,\ldots,m_r),
 \qquad m_i>0
\end{equation*}
after omitting zero singular values. The kinetic terms are unchanged because the
changes are unitary. For each nonzero diagonal pair there remains the symmetry
\begin{equation*}
 \nu_i(q)\mapsto\ee^{\ii\alpha_i}\nu_i(q),
 \qquad
 \nu_i(-q)\mapsto\ee^{-\ii\alpha_i}\nu_i(-q),
\end{equation*}
which is the fermion-number $\mathrm U(1)$ of a Dirac field. These are the
approximate flavor symmetries when weak interactions are neglected. In the basis
with real positive masses, the standard $C$, $P$, and $T$ actions can be chosen to
preserve the mass term; their phases are absorbed into the two Weyl fields.

For $q=0$, the coefficient is a complex symmetric matrix $M=M^T$. Its canonical
form is the Takagi factorization
\begin{equation*}
 M=V^*D V^\dagger,
 \qquad V^T M V=D,
 \qquad D=\operatorname{diag}(m_1,\ldots,m_n),\quad m_i\geq0.
\end{equation*}
There is one unitary change of basis, rather than two independent changes. After
the Takagi transformation the continuous phase symmetry of an isolated massive
field is absent; only its fermion-parity sign remains, with extra orthogonal
rotations possible inside degenerate mass subspaces.

Under CP a Weyl multiplet transforms as
$\nu\mapsto X\,\ii\sigma_2\nu^*$, and invariance of the Majorana term requires
\begin{equation*}
 X^T M X=M^*.
\end{equation*}
Time reversal gives the analogous condition with the anti-unitary complex
conjugation. In the Takagi basis $M=D$ is real, so canonical $C$, CP, and $T$
actions can be chosen whenever the internal representation itself admits the
corresponding intertwiner. A generic charged Dirac phase has been removed by the
two-field freedom, while a Majorana phase is physical up to the residual discrete
signs and degeneracy rotations.

% BANKS-SOURCE: pdf=70 print=60 kind=solution id=5.6
\BanksSolution{5.6}
Insert a complete set of momentum eigenstates between the fields and define the
positive-frequency spectral matrix. Write $s=-p^2\geq0$ for its invariant mass
variable.
\begin{equation*}
 \rho_{\alpha\beta}(p)=\sum_n(2\pi)^D\delta^D(p-p_n)
 \langle0|\psi_\alpha(0)|n\rangle
 \langle n|\bar\psi_\beta(0)|0\rangle.
\end{equation*}
Translation invariance makes the two-point function a Fourier transform of this
measure. Lorentz covariance and parity allow only two Dirac structures,
\begin{equation*}
 \rho(p)=-\rho_1(s)\slashed p-\rho_2(s),
 \qquad \rho_i(s)=0\quad(s<0).
\end{equation*}
Positivity of the state sum implies $\rho_1(s)\geq0$ and, on the support of
the measure, $\sqrt{s}\,\rho_1(s)\geq|\rho_2(s)|$. The spectral representation of
the Feynman propagator is therefore
\begin{equation*}
 S_F(p)=\ii\int_0^\infty\dd s\,
 \frac{-\rho_1(s)\slashed p-\rho_2(s)}{p^2+s-\ii\epsilon}
 +\text{polynomial contact terms}.
\end{equation*}
For a stable one-particle pole of mass $m$,
$\rho_1(s)=Z\delta(s-m^2)$ and $\rho_2(s)=Zm\delta(s-m^2)$, giving
$S_F(p)=-\ii Z(\slashed p+m)/(p^2+m^2-\ii\epsilon)$ at the pole. If parity is
not imposed, the two allowed structures $\slashed p\gamma_5$ and $\gamma_5$
must be added with their own spectral measures.

For a conserved current define
\begin{equation*}
 \rho_{\mu\nu}(p)=\sum_n(2\pi)^D\delta^D(p-p_n)
 \langle0|J_\mu(0)|n\rangle\langle n|J_\nu(0)|0\rangle.
\end{equation*}
Current conservation gives $p^\mu\rho_{\mu\nu}=0$, and Lorentz covariance then
forces
\begin{equation*}
 \rho_{\mu\nu}(p)=\left(p_\mu p_\nu-p^2\eta_{\mu\nu}\right)\rho_J(p^2),
 \qquad \rho_J(s)\geq0.
\end{equation*}
Consequently,
\begin{equation*}
 \Pi_{\mu\nu}(p)
 =-\ii\int_0^\infty\dd s\,
 \frac{p_\mu p_\nu+s\eta_{\mu\nu}}{p^2+s-\ii\epsilon}\rho_J(s)
 +\text{local subtraction polynomial}.
\end{equation*}
Factoring the transverse tensor defines the scalar function
$\Pi(p^2)$, so that
\begin{equation*}
 \Pi_{\mu\nu}(p)
 =(p_\mu p_\nu-p^2\eta_{\mu\nu})\Pi(p^2).
\end{equation*}
The subtraction polynomial is allowed by short-distance singularities and is
itself transverse after contact terms are combined. A massless vector state appears
as a $\delta(s)$ contribution, with the transverse expression understood as a
distribution at $s=0$.

% BANKS-SOURCE: pdf=70 print=60 kind=solution id=5.7
\BanksSolution{5.7}
Define the covariant momentum and field strength by
\begin{equation*}
 \Pi_\mu=-\ii\partial_\mu-eA_\mu,
 \qquad
 \mathcal F_{\mu\nu}=\partial_\mu A_\nu-\partial_\nu A_\mu.
\end{equation*}
For the potential in the problem, $\mathcal F$ is constant (if $F$ is taken
antisymmetric, $\mathcal F_{\mu\nu}=-2F_{\mu\nu}$). Since
$[\Pi_\mu,\Pi_\nu]=+\ii e\mathcal F_{\mu\nu}$,
\begin{equation*}
 (\slashed\Pi)^2=-\Pi^2+\frac e2\sigma^{\mu\nu}\mathcal F_{\mu\nu},
 \qquad \sigma^{\mu\nu}=\frac\ii2[\gamma^\mu,\gamma^\nu].
\end{equation*}
Thus a Green function satisfying the equation in the problem is
\begin{equation*}
 S(x,y)=-(\slashed\Pi_x+m)G(x,y),
 \qquad
 G=\left[\Pi^2-\frac e2\sigma^{\mu\nu}\mathcal F_{\mu\nu}+m^2-\ii\epsilon\right]^{-1}_{x,y}.
\end{equation*}
The inverse has the Schwinger representation
\begin{equation*}
 G(x,y)=+\ii\int_0^\infty\dd s\,
 \ee^{-\ii s(m^2-\ii\epsilon)}
 \left\langle x\left|\ee^{-\ii s[\Pi^2-(e/2)\sigma\cdot\mathcal F]}\right|y\right\rangle.
\end{equation*}
To display the kernel, put $z=x-y$, use the straight Wilson factor
\begin{equation*}
 \Phi(x,y)=\exp\!\left[+\ii e\int_y^x A_\mu(\xi)\,\dd\xi^\mu\right],
\end{equation*}
and regard $e\mathcal F$ as a matrix on Lorentz indices. The constant-field
heat kernel is
\begin{align*}
 K(x,y;s)={}&\Phi(x,y)\frac{1}{(4\pi\ii s)^{D/2}}
 \det{}^{-1/2}\!\left[\frac{\sin(e\mathcal F s)}{e\mathcal F s}\right]\\
 &\times\exp\!\left[+\frac\ii4 z_\mu
   (e\mathcal F\cot(e\mathcal F s))^{\mu}{}_{\nu}z^\nu\right]
 \exp\!\left[+\frac{\ii e s}{2}\sigma^{\mu\nu}\mathcal F_{\mu\nu}\right].
\end{align*}
Therefore $G=+\ii\int_0^\infty\dd s\,\ee^{-\ii s(m^2-\ii\epsilon)}K$ and
$S=-(\slashed\Pi_x+m)G$. The matrix functions are defined by their power
series, or by continuation from Euclidean signature. As $\mathcal F\to0$,
$\sin(e\mathcal F s)/(e\mathcal F s)\to1$ and
$e\mathcal F\cot(e\mathcal F s)\to s^{-1}$, so the free Dirac Green function
is recovered. A convention in which the time-ordered propagator is
$\ii S$ only multiplies this solution by $\ii$.

% BANKS-SOURCE: pdf=70 print=60 kind=solution id=5.8
\BanksSolution{5.8}
In $D=4$, in the Weyl basis,
\begin{equation*}
 \gamma^0=\begin{pmatrix}0&1\\1&0\end{pmatrix},
 \qquad
 \gamma^i=\begin{pmatrix}0&\sigma_i\\-\sigma_i&0\end{pmatrix}.
\end{equation*}
The Pauli identity
$\sigma_i\sigma_j=\delta_{ij}1+\ii\epsilon_{ijk}\sigma_k$ gives
\begin{equation*}
 \sigma^\mu\bar\sigma^\nu+\sigma^\nu\bar\sigma^\mu
 =-2\eta^{\mu\nu}1.
\end{equation*}
Multiplying the block matrices proves
\begin{equation*}
 [\gamma^\mu,\gamma^\nu]_+=-2\eta^{\mu\nu}1.
\end{equation*}
The ordered product of the four matrices is block diagonal, and direct use of
the Pauli matrices gives
\begin{equation*}
 \gamma_5\equiv\ii\gamma^0\gamma^1\gamma^2\gamma^3
 =\begin{pmatrix}-1&0\\0&1\end{pmatrix}.
\end{equation*}
To square it, move the second copy of each $\gamma^\mu$ through the first copy:
\begin{equation*}
 (\gamma^0\gamma^1\gamma^2\gamma^3)^2
 =(-1)^6(\gamma^0)^2(\gamma^1)^2(\gamma^2)^2(\gamma^3)^2=-1,
\end{equation*}
because $(\gamma^0)^2=1$ and $(\gamma^i)^2=-1$. The prefactor $\ii^2=-1$
then gives $\gamma_5^2=1$.

Finally, moving a given $\gamma^\mu$ through the other three factors changes
the sign three times. Thus
\begin{equation*}
 \gamma_5\gamma^\mu=-\gamma^\mu\gamma_5,
 \qquad [\gamma_5,\gamma^\mu]_+=0.
\end{equation*}
The same conclusion follows immediately from the displayed block form.

% BANKS-SOURCE: pdf=70 print=60 kind=solution id=5.9
\BanksSolution{5.9}
Let $\chi_i,\bar\chi_i$ be independent Grassmann variables and take the
Gaussian weight $\exp(-\bar\chi A\chi)$, with $A$ invertible. Introduce Grassmann
sources $\bar\eta,\eta$ and shift the variables:
\begin{align*}
 Z[\bar\eta,\eta]
 &=\int[d\bar\chi\,d\chi]\,
   \exp\left[-\bar\chi A\chi+\bar\eta\chi+\bar\chi\eta\right]\\
 &=\det A\,\exp(\bar\eta A^{-1}\eta).
\end{align*}
Translation invariance of the Berezin integral justifies the shift
$\chi\mapsto\chi+A^{-1}\eta$ and
$\bar\chi\mapsto\bar\chi+\bar\eta A^{-1}$. Differentiating with ordered left
derivatives generates the normalized correlation functions. Since the exponent
is bilinear in the sources, an odd number of fields gives zero and
\begin{equation*}
 \left\langle\chi_{i_1}\cdots\chi_{i_n}
 \bar\chi_{j_1}\cdots\bar\chi_{j_n}\right\rangle
 =\det\left[(A^{-1})_{i_rj_s}\right]_{r,s=1}^n
 =\sum_{\pi\in S_n}\operatorname{sgn}(\pi)
   \prod_{r=1}^n(A^{-1})_{i_rj_{\pi(r)}}.
\end{equation*}
The sign is the sign of the permutation needed to put each contracted pair next
to one another. For an arbitrary ordered list $X_1,\ldots,X_{2n}$ of fields,
the same statement is
\begin{equation*}
 \langle X_1\cdots X_{2n}\rangle
 =\sum_{\mathcal P}\operatorname{sgn}(\mathcal P)
   \prod_{(r,s)\in\mathcal P}\langle X_rX_s\rangle,
\end{equation*}
where the sum is over pairings and the sign is the parity of the crossings of
the pairs. This is fermionic Wick's theorem. For example,
\begin{equation*}
 \langle X_1X_2X_3X_4\rangle
 =\langle X_1X_2\rangle\langle X_3X_4\rangle
 -\langle X_1X_3\rangle\langle X_2X_4\rangle
 +\langle X_1X_4\rangle\langle X_2X_3\rangle.
\end{equation*}
The minus sign in the middle term is the exchange sign that distinguishes the
Grassmann theorem from a commuting Gaussian.

% BANKS-SOURCE: pdf=70 print=60 kind=solution id=5.10
\BanksSolution{5.10}
For an antisymmetric $N\times N$ matrix $A$, write
$\bar\eta=X_++X_-$ and $\eta=X_+-X_-$, where
$X_\pm=(\bar\eta\pm\eta)/2$. Grassmann transposition and $A^T=-A$ give
\begin{equation*}
 \bar\eta^TA\eta=X_+^TAX_+-X_-^TAX_-.
\end{equation*}
The mixed terms cancel because
$X_-^TAX_+=X_+^TAX_-$. Put $N=2r$. In block ordering, the change of
variables is
\begin{equation*}
 \begin{pmatrix}\bar\eta\\ \eta\end{pmatrix}
 =
 \underbrace{\begin{pmatrix}1&1\\1&-1\end{pmatrix}}_{B}
 \begin{pmatrix}X_+\\X_-\end{pmatrix},
 \qquad \det B=(-2)^N.
\end{equation*}
For Grassmann variables the measure transforms with the inverse determinant.
The permutation from the original pairwise ordering to block ordering contributes
$(-1)^{N(N-1)/2}=(-1)^r$, so the determinant orientation is
\begin{equation*}
 [d^N\bar\eta\,d^N\eta]
 =(-1)^r2^{-N}[d^NX_+\,d^NX_-].
\end{equation*}
This records the Berezinian and the orientation rather than hiding their sign.
Using $X^TAX=\frac12X^T(2A)X$, the determinant integral becomes
\begin{align*}
 \int[d\bar\eta\,d\eta]\,\ee^{\bar\eta^TA\eta}
 &=(-1)^r2^{-N}\operatorname{Pf}(2A)\operatorname{Pf}(-2A)\\
 &=(-1)^r2^{-2r}\left(2^r\operatorname{Pf}(A)\right)
       \left((-1)^r2^r\operatorname{Pf}(A)\right)
 =\operatorname{Pf}(A)^2.
\end{align*}
For odd $N$, both the Pfaffian integral and the determinant integral vanish.

For completeness, expansion of the exponential gives
\begin{equation*}
 \operatorname{Pf}(A)=\frac1{2^r r!}
 \epsilon_{i_1\ldots i_{2r}}
 A_{i_1i_2}\cdots A_{i_{2r-1}i_{2r}}.
\end{equation*}
The remaining sign and normalization can be checked without integration. Under a
congruence $A\mapsto S^TAS$, the determinant acquires $(\det S)^2$, while the
Pfaffian acquires $\det S$. Every complex antisymmetric matrix is congruent to a
direct sum of blocks
\begin{equation*}
 J(a_j)=\begin{pmatrix}0&a_j\\-a_j&0\end{pmatrix}
\end{equation*}
and zero blocks. For one block, $\det J(a)=a^2$ and
$\operatorname{Pf}J(a)=a$ with the stated orientation. Multiplying the blocks
therefore proves, for every antisymmetric matrix,
\begin{equation*}
 \boxed{\det A=\operatorname{Pf}(A)^2}.
\end{equation*}

% BANKS-SOURCE: pdf=70 print=60 kind=solution id=5.11
\BanksSolution{5.11}
In $D=4$, take the interaction-picture $S$ matrix to be the time-ordered exponential
of $\ii\int\dd^D x\,\Delta\mathcal L$. A single insertion contributes
\begin{equation*}
 \ii\phi\left(g_S\bar\psi\psi+g_P\bar\psi\gamma_5\psi\right),
\end{equation*}
so the Yukawa vertex, with the fermion-number arrow entering through the
$\psi$ leg and leaving through the $\bar\psi$ leg, is
\begin{equation*}
 \boxed{\quad \ii\left(g_S+g_P\gamma_5\right)\quad}.
\end{equation*}
The scalar line contributes
$-\ii/(k^2+m_\phi^2-\ii\epsilon)$, and a free fermion line oriented in the
direction of fermion number contributes
\begin{equation*}
 \frac{\ii}{\slashed p-m+\ii\epsilon}
 =-\frac{\ii(\slashed p+m)}{p^2+m^2-\ii\epsilon}.
\end{equation*}
At every vertex impose four-momentum conservation. Along an open fermion line
keep the matrices in arrow order. For example, a tree-level exchange between
two external fermions contains
\begin{equation*}
 \bar u(p')\,\ii(g_S+g_P\gamma_5)u(p)
 \frac{-\ii}{(p-p')^2+m_\phi^2-\ii\epsilon}
 \bar u(k')\,\ii(g_S+g_P\gamma_5)u(k),
\end{equation*}
with the corresponding $u\leftrightarrow v$ substitutions for external
antiparticles and the reversed arrow for an antiparticle line. External incoming
fermions use $u$, outgoing fermions use $\bar u$, incoming antifermions use
$\bar v$, and outgoing antifermions use $v$.

Hermiticity fixes the allowed coupling convention. Since
$(\bar\psi\psi)^\dagger=\bar\psi\psi$ and
$(\bar\psi\gamma_5\psi)^\dagger=-\bar\psi\gamma_5\psi$, real $g_S$ and
purely imaginary $g_P$ make the displayed interaction Hermitian. If one instead
writes the Hermitian pseudoscalar coupling as
$\ii g_5\phi\bar\psi\gamma_5\psi$ with real $g_5$, its vertex is
$-g_5\gamma_5$. The scalar and fermion propagators and the arrow rules are
unchanged.

% BANKS-SOURCE: pdf=70 print=60 kind=solution id=5.12
\BanksSolution{5.12}
In $D$ dimensions the Clifford relation first gives
$\gamma^\lambda\gamma_\lambda=-D$ and
$\gamma^\lambda\gamma_\mu=-\gamma_\mu\gamma^\lambda-2\delta^\lambda_\mu$.
Moving the leftmost matrix through a string proves
\begin{equation*}
 \gamma^\lambda\gamma_\mu\gamma_\lambda=(D-2)\gamma_\mu,
 \qquad
 \gamma^\lambda\gamma_\mu\gamma_\nu\gamma_\lambda=4\eta_{\mu\nu}\quad(D=4).
\end{equation*}
For antisymmetrized products the induction takes the compact
form
\begin{equation*}
 \gamma^\lambda\gamma_{\mu_1\ldots\mu_r}\gamma_\lambda
 =(-1)^{r+1}(D-2r)\gamma_{\mu_1\ldots\mu_r}.
\end{equation*}
In $D=4$ it includes the next cases
\begin{equation*}
 \gamma^\lambda\gamma_{\mu\nu\rho}\gamma_\lambda=-2\gamma_{\mu\nu\rho},
 \qquad
 \gamma^\lambda\gamma_{\mu\nu\rho\sigma}\gamma_\lambda=4\gamma_{\mu\nu\rho\sigma}.
\end{equation*}
For an ordinary product, expand it into antisymmetric grades and apply the same
formula term by term. The rank-four case displayed above is already a higher
power than the identities stated in Appendix C.

Cyclicity, anticommutation with $\gamma_5$, and a single insertion of
$\eta_{\mu\nu}$ determine all ordinary traces recursively in $D=4$:
\begin{align*}
 \Tr 1&=4, & \Tr(\gamma^\mu)&=0,\\
 \Tr(\gamma^\mu\gamma^\nu)&=-4\eta^{\mu\nu},\\
 \Tr(\gamma^\mu\gamma^\nu\gamma^\rho\gamma^\sigma)
 &=4\left(\eta^{\mu\nu}\eta^{\rho\sigma}
 -\eta^{\mu\rho}\eta^{\nu\sigma}
 +\eta^{\mu\sigma}\eta^{\nu\rho}\right).
\end{align*}
All odd traces vanish: moving $\gamma_5$ cyclically through an odd product
changes the sign while leaving the trace unchanged. More generally,
\begin{equation*}
 \Tr(\gamma^{\mu_1}\cdots\gamma^{\mu_{2r}})
 =\sum_{j=2}^{2r}(-1)^{j-1}\eta^{\mu_1\mu_j}
 \Tr(\gamma^{\mu_2}\cdots\widehat{\gamma^{\mu_j}}\cdots\gamma^{\mu_{2r}}),
\end{equation*}
with the four-gamma result as the first nontrivial seed.

With $\epsilon^{0123}=+1$ and the definition of $\gamma_5$ in $D=4$,
\begin{align*}
 \Tr\gamma_5&=0, & \Tr(\gamma_5\gamma^{\mu_1}\cdots\gamma^{\mu_r})&=0
       &&(r=1,2,3),\\
 \Tr(\gamma_5\gamma^\mu\gamma^\nu\gamma^\rho\gamma^\sigma)
 &=-4\ii\epsilon^{\mu\nu\rho\sigma}.
\end{align*}
For any even number of matrices, the higher traces are obtained by contracting
down to the four-gamma seed. In particular, the six-gamma case is the signed sum
of the fifteen terms
\begin{equation*}
 \Tr(\gamma_5\gamma^{\mu_1}\cdots\gamma^{\mu_6})
 =-4\ii\sum_{\substack{\text{one pair }(i,j)\\\text{and ordered remainder }a,b,c,d}}
 \operatorname{sgn}(i,j;a,b,c,d)\,
 \eta^{\mu_i\mu_j}\epsilon^{\mu_a\mu_b\mu_c\mu_d}.
\end{equation*}
The sign is the parity of the permutation taking
$(1,2,3,4,5,6)$ to $(i,j,a,b,c,d)$. Finally, inserting $\gamma_5$ into a
contraction gives
\begin{equation*}
 \gamma^\lambda\gamma_{\mu_1\ldots\mu_r}\gamma_5\gamma_\lambda
 =(-1)^r(4-2r)\gamma_{\mu_1\ldots\mu_r}\gamma_5,
\end{equation*}
which completes the requested trace and contraction identities.

% BANKS-SOURCE: pdf=70 print=60 kind=solution id=5.13
\BanksSolution{5.13}
The auxiliary fields occur algebraically. Varying with respect to $\phi^A$ gives
\begin{equation*}
 \phi^A+g_A\bar\psi\Gamma^A\psi=0,
 \qquad
 \phi^A=-g_A\bar\psi\Gamma^A\psi.
\end{equation*}
Substitution, or completion of the square, gives
\begin{equation*}
 \mathcal L_{\rm aux}
 =\frac12\left(\phi^A+g_A\bar\psi\Gamma^A\psi\right)^2
 -\frac12g_A^2(\bar\psi\Gamma^A\psi)(\bar\psi\Gamma_A\psi),
\end{equation*}
so integrating out $\phi$ produces the four-fermion term
\begin{equation*}
 \mathcal L_4=-\frac12g_A^2
 (\bar\psi\Gamma^A\psi)(\bar\psi\Gamma_A\psi).
\end{equation*}
In $D=4$, the sixteen Dirac matrices decompose into the five Lorentz channels
$1$, $\gamma_5$, $\gamma^\mu$, $\gamma^\mu\gamma_5$, and
$\gamma^{\mu\nu}$. Giving each irreducible channel its own coefficient and
contracting its Lorentz indices supplies the scalar, pseudoscalar, vector,
axial-vector, and tensor four-fermion terms. Fierz identities relate different
ways of pairing the four spinors, so this channel basis spans the most general
Lorentz-invariant interaction quadratic in each of $\psi$ and $\bar\psi$.
Every bilinear $\bar\psi\Gamma^A\psi$ is invariant under
$\psi\mapsto\ee^{\ii\alpha}\psi$, so the resulting interaction preserves the
global $\mathrm U(1)$ symmetry. Auxiliary fields with other tensor structures or
higher powers give the corresponding generalizations, including terms that do
not preserve this phase symmetry.

For the functional integral, write
$D[\phi]=-\ii\slashed\partial-\phi_A\Gamma^A$ and $D=D_0-V$. Grassmann
integration gives
\begin{equation*}
 \int[d\bar\psi\,d\psi]\,\ee^{\ii\int\bar\psi D\psi}
 =\operatorname{Det}D,
 \qquad
 \log\operatorname{Det}D
 =\log\operatorname{Det}D_0
 -\sum_{n\geq1}\frac1n\Tr\left[(D_0^{-1}V)^n\right].
\end{equation*}
Each functional trace closes the directed fermion line and is therefore one
closed fermion loop. A commuting bosonic Gaussian instead gives
$\operatorname{Det}^{-1}D$ and reverses the sign of every trace in
$\log\operatorname{Det}D$. Hence every closed fermion loop contributes one extra
factor of $-1$, independent of the number of vertices on that loop.

% BANKS-SOURCE: pdf=71 print=61 kind=solution id=5.14
\BanksSolution{5.14}
Use the source convention
$\gamma^{\mu\nu}=\frac12(\gamma^\nu\gamma^\mu-\gamma^\mu\gamma^\nu)$,
which is the same convention used in the displayed identity. With the
mostly-plus Clifford relation, the on-shell equations imply
\begin{equation*}
 \bar u(k)\slashed k=m\bar u(k),
 \qquad \slashed p\,u(p)=m u(p).
\end{equation*}
Since
\begin{equation*}
 \gamma^{\mu\nu}(k-p)_\nu
 =\frac12\left[(\slashed k-\slashed p)\gamma^\mu
 -\gamma^\mu(\slashed k-\slashed p)\right],
\end{equation*}
the anticommutator gives, between the spinors,
\begin{align*}
 \bar u(k)\slashed k\gamma^\mu u(p)&=m\bar u(k)\gamma^\mu u(p),\\
 \bar u(k)\slashed p\gamma^\mu u(p)&=-2p^\mu\bar u(k)u(p)-m\bar u(k)\gamma^\mu u(p),\\
 \bar u(k)\gamma^\mu\slashed k u(p)&=-2k^\mu\bar u(k)u(p)-m\bar u(k)\gamma^\mu u(p),\\
 \bar u(k)\gamma^\mu\slashed p u(p)&=m\bar u(k)\gamma^\mu u(p).
\end{align*}
Combining the four lines yields
\begin{equation*}
 \bar u(k)\gamma^{\mu\nu}(k-p)_\nu u(p)
 =2m\bar u(k)\gamma^\mu u(p)+(p+k)^\mu\bar u(k)u(p),
\end{equation*}
which is exactly
\begin{equation*}
 2m\bar u(k)\gamma^\mu u(p)=\bar u(k)
 \left[-(p+k)^\mu+\gamma^{\mu\nu}(k-p)_\nu\right]u(p).
\end{equation*}

For antiparticle spinors, $\bar v(k)\slashed k=-m\bar v(k)$ and
$\slashed p\,v(p)=-m v(p)$. The same calculation gives
\begin{equation*}
 2m\bar v(k)\gamma^\mu v(p)=\bar v(k)
 \left[(p+k)^\mu-\gamma^{\mu\nu}(k-p)_\nu\right]v(p).
\end{equation*}
Equivalently, on reversing the labels,
\begin{equation*}
 2m\bar v(p)\gamma^\mu v(k)=\bar v(p)
 \left[(p+k)^\mu-\gamma^{\mu\nu}(p-k)_\nu\right]v(k).
\end{equation*}
The relative signs record the Clifford relation and the negative mass eigenvalue
in the $v$ equation.

% BANKS-SOURCE: pdf=71 print=61 kind=solution id=5.15
\BanksSolution{5.15}
In $D=4$, let $E=\omega_{\mathbf p}=(\mathbf p^2+m^2)^{1/2}$ and choose orthonormal two-spinors
$\chi_s$. In the Dirac basis,
\begin{equation*}
 \gamma_D^0=\begin{pmatrix}1&0\\0&-1\end{pmatrix},
 \qquad
 \gamma_D^i=\begin{pmatrix}0&\sigma_i\\-\sigma_i&0\end{pmatrix}.
\end{equation*}
Writing $u=(\varphi,\xi)^T$, the equation
$(\slashed p-m)u=0$ gives
$-(E+m)\varphi+(\boldsymbol\sigma\cdot\mathbf p)\xi=0$ and
$-(\boldsymbol\sigma\cdot\mathbf p)\varphi+(E-m)\xi=0$. A convenient normalized
solution is
\begin{equation*}
 u_D(p,s)=\begin{pmatrix}
  \dfrac{\boldsymbol\sigma\cdot\mathbf p}{\sqrt{E+m}}\,\chi_s\\[2pt]
  \sqrt{E+m}\,\chi_s
 \end{pmatrix}.
\end{equation*}
For $(\slashed p+m)v=0$, choose an orthonormal spin basis $\eta_s$ and obtain
\begin{equation*}
 v_D(p,s)=\begin{pmatrix}
  \sqrt{E+m}\,\eta_s\\[2pt]
  \dfrac{\boldsymbol\sigma\cdot\mathbf p}{\sqrt{E+m}}\,\eta_s
 \end{pmatrix}.
\end{equation*}
One may take $\eta_s=-\ii\sigma_2\chi_s^*$; other choices are spin-basis
phases. These spinors obey
\begin{equation*}
  \bar u(p,r)u(p,s)=-2m\delta_{rs},
  \qquad \bar v(p,r)v(p,s)=2m\delta_{rs},
\end{equation*}
and the completeness relations in Appendix C.

In the Weyl basis use
\begin{equation*}
 \gamma_W^\mu=\begin{pmatrix}0&\sigma^\mu\\\bar\sigma^\mu&0\end{pmatrix},
  \qquad
  -p_\mu\sigma^\mu=E-\boldsymbol\sigma\cdot\mathbf p,
  \qquad
  -p_\mu\bar\sigma^\mu=E+\boldsymbol\sigma\cdot\mathbf p,
\end{equation*}
with $p_\mu=(-E,\mathbf p)$. Since
$(-p_\mu\sigma^\mu)(-p_\nu\bar\sigma^\nu)=m^2$, their positive Hermitian square roots
exist. The two independent solutions are
\begin{equation*}
 u_W(p,s)=\begin{pmatrix}
  \sqrt{-p_\mu\sigma^\mu}\,\chi_s\\[2pt]
  -\sqrt{-p_\mu\bar\sigma^\mu}\,\chi_s
 \end{pmatrix},
 \qquad
 v_W(p,s)=\begin{pmatrix}
  \sqrt{-p_\mu\sigma^\mu}\,\eta_s\\[2pt]
  \sqrt{-p_\mu\bar\sigma^\mu}\,\eta_s
 \end{pmatrix}.
\end{equation*}
Substitution uses the two matrix products above and proves the momentum-space
Dirac equations. Applying $S_D^\dagger$ to these Weyl spinors gives the displayed
Dirac-basis spinors, up to the allowed unitary change of the two-spinor basis.

% BANKS-SOURCE: pdf=71 print=61 kind=solution id=5.16
\BanksSolution{5.16}
Let $R(g)$ be the unitary matrix representing $g\in G$ on a multiplet of fields.
Charge conjugation maps a field into its complex conjugate, so invariance of the
representation means that, choosing the intertwiner in the convenient direction,
there is a unitary matrix $C$ such that
\begin{equation*}
 C R(g)^* C^{-1}=R(g).
\end{equation*}
Equivalently, the anti-linear map $J=CK$, where $K$ is componentwise complex
conjugation, commutes with every $R(g)$:
\begin{equation*}
 J R(g)=R(g)J.
\end{equation*}
On an irreducible self-conjugate representation, Schur's lemma says that
$J^2$ is a scalar, say $J^2=\lambda1$. Anti-unitarity gives $|\lambda|=1$,
while $J J^2=J^2J$ gives $\lambda^*=\lambda$. Thus
in matrix language,
\begin{equation*}
 C C^*=J^2=+1\quad\text{or}\quad C C^*=J^2=-1.
\end{equation*}
The phase of an anti-linear map does not change its square. Unitarity of $C$ and
$CC^*=\lambda1$ imply $C^T=\lambda C$. Define the bilinear form
$B(v,w)=\langle Jv,w\rangle$, whose matrix is $C^\dagger$ in the standard
inner product. The intertwining relation gives
\begin{equation*}
 B(R(g)v,R(g)w)=B(v,w),
 \qquad B^T=\lambda B,
\end{equation*}
A symmetric invariant bilinear form, or $J^2=+1$, is the real case. An
antisymmetric invariant form, or $J^2=-1$, is the pseudo-real (quaternionic)
case. If an irreducible representation is complex, it has no such $J$ by itself;
it must occur together with its conjugate, which charge conjugation exchanges.
Thus each self-conjugate irreducible allowed by charge conjugation is real or
pseudo-real, while complex irreducibles occur in conjugate pairs.

\end{bankssolutions}
